\section{Alignment Data}
\label{sec:data}

We define the \textbf{Superficial Alignment Hypothesis}:
% \begin{quote}
A model's knowledge and capabilities are learnt almost entirely during pretraining, while alignment teaches it which subdistribution of formats should be used when interacting with users.
% \end{quote}
If this hypothesis is correct, and alignment is largely about learning style, then a corollary of the \hypothesis{} is that one could sufficiently tune a pretrained language model with a rather small set of examples \citep{kirstain2021examples}.

To that end, we collect a dataset of 1,000 prompts and responses, where the outputs (responses) are stylistically aligned with each other, but the inputs (prompts) are diverse.
Specifically, we seek outputs in the style of a helpful AI assistant.
We curate such examples from a variety of sources, primarily split into community Q\&A forums and manually authored examples.
We also collect a test set of 300 prompts and a development set of 50.
Table~\ref{tab:data_sources} shows an overview of the different data sources and provides some statistics (see Appendix~\ref{sec:training_examples} for a selection of training examples).%\footnote{We highly encourage the reader to review the sample of training examples in Appendix~\ref{sec:training_examples}.}

\begin{table}[t]
    \small
    \centering
    \begin{tabular}{@{}lrrr@{}}
        \toprule
        \textbf{Source} & \textbf{\#Examples} & \textbf{Avg Input Len.} & \textbf{Avg Output Len.} \\
        \midrule
        \textbf{Training} & & & \\
        \quad Stack Exchange (STEM) & 200 & 117 & 523\\
        \quad Stack Exchange (Other) & 200 & 119 & 530 \\
        \quad wikiHow & 200 & 12 & 1,811 \\
        \quad Pushshift r/WritingPrompts & 150 & 34 & 274 \\
        \quad Natural Instructions & 50 & 236 & 92\\
        \quad Paper Authors (Group A) & 200 & 40 & 334 \\
        \midrule
        \textbf{Dev} & & & \\
        \quad Paper Authors (Group A) & 50 & 36 & N/A \\
        \midrule
        \textbf{Test} & & & \\
        \quad Pushshift r/AskReddit & 70 & 30 & N/A \\
        \quad Paper Authors (Group B) & 230 & 31 & N/A \\
        \bottomrule
    \end{tabular}
    \caption{Sources of training prompts (inputs) and responses (outputs), and test prompts. The total amount of training data is roughly 750,000 tokens, split over exactly 1,000 sequences.}
    \label{tab:data_sources}
\end{table}


\subsection{Community Questions \& Answers}
\label{sec:community_data}

We collect data from three community Q\&A websites: Stack Exchange, wikiHow, and the Pushshift Reddit Dataset \citep{baumgartner2020pushshift}.
% We collect data from three community Q\&A websites: Stack Exchange,\footnote{\url{https://stackexchange.com/}} wikiHow,\footnote{\url{https://www.wikihow.com/}} and the Pushshift Reddit Dataset \citep{baumgartner2020pushshift}.
Largely speaking, answers from Stack Exchange and wikiHow are well-aligned with the behavior of a helpful AI agent, and can therefore be mined automatically, whereas highly upvoted Reddit answers tend to be humorous or trolling, requiring a more manual approach to curate responses that follow the appropriate style.

\paragraph{Stack Exchange}
Stack Exchange contains 179 online communities (exchanges), each one dedicated to a specific topic, with the most popular one being programming (Stack Overflow).
Users can post questions, answers, comments and upvote (or downvote) all of the above.
Thanks to active community members and moderators, Stack Exchange has successfully maintained a high bar for content quality.

We apply both quality and diversity controls when sampling from Stack Exchange.
First, we divide the exchanges into 75 STEM exchanges (including programming, math, physics, etc.) and 99 other (English, cooking, travel, and more); we discard 5 niche exchanges.
We then sample 200 questions and answers from each set using a temperature of $\tau = 3$ to get a more uniform sample of the different domains.
Within each exchange, we take the questions with the highest score that are self-contained in the title (no body).
We then select the top answer for each question, assuming it had a strong positive score (at least 10).
To conform with the style of a helpful AI assistant, we automatically filter answers that are too short (less than 1200 characters), too long (more than 4096 characters), written in the first person (``\texttt{I}'', ``\texttt{my}''), or reference other answers (``\texttt{as mentioned}'', ``\texttt{stack exchange}'', etc);
we also remove links, images, and other HTML tags from the response, retaining only code blocks and lists.
Since Stack Exchange questions contain both a title and a description, we randomly select the title as the prompt for some examples, and the description for others.
% Figure~\ref{fig:training_examples} shows two of the selected examples from Stack Exchange.



\paragraph{wikiHow}
wikiHow is an online wiki-style publication featuring over 240,000 how-to articles on a variety of topics.
Anyone can contribute to wikiHow, though articles are heavily moderated, resulting in almost universally high-quality content.
We sample 200 articles from wikiHow, sampling a category first (out of 19) and then an article within it to ensure diversity.
We use the title as the prompt (e.g. ``\texttt{How to cook an omelette?}'') and the article's body as the response.
We replace the typical ``\texttt{This article...}'' beginning with ``\texttt{The following answer...}'', and apply a number of preprocessing heuristics to prune links, images, and certain sections of the text.
% Figure~\ref{fig:training_examples} shows one of the selected examples from wikiHow.

\paragraph{The Pushshift Reddit Dataset}
Reddit is one of the most popular websites in the world, allowing users to share, discuss, and upvote content in user-created subreddits.
Due to its immense popularity, Reddit is geared more towards entertaining fellow users rather than helping; it is quite often the case that witty, sarcastic comments will obtain more votes than serious, informative comments to a post.
We thus restrict our sample to two subsets, r/AskReddit and r/WritingPrompts, and manually select examples from within the most upvoted posts in each community.
From r/AskReddit we find 70 self-contained prompts (title only, no body), which we use for the test set, since the top answers are not necessarily reliable.
The WritingPrompts subreddit contains premises of fictional stories, which other users are then encouraged to creatively complete.
We find 150 prompts and high-quality responses, encompassing topics such as love poems and short science fiction stories, which we add to the training set.
All data instances were mined from the Pushshift Reddit Dataset \citep{baumgartner2020pushshift}.%\footnote{\url{https://pushshift.io/}}


\subsection{Manually Authored Examples}
\label{sec:manual_data}

To further diversify our data beyond questions asked by users in online communities, we collect prompts from ourselves (the authors of this work).
We designate two sets of authors, Group A and Group B, to create 250 prompts each, inspired by their own interests or those of their friends.\footnote{Despite our efforts to prevent leakage, there was significant contact between the groups before the annotation process, which resulted in certain shared priors that can be observed in the data.}
We select 200 prompts from Group A for training and 50 prompts as a held-out development set.
After filtering some problematic prompts, the remaining 230 prompts from Group B are used for test.

We supplement the 200 training prompts with high-quality answers, which we write ourselves.
While authoring answers, we try to set a uniform tone that is appropriate for a helpful AI assistant.
Specifically, many prompts will be answered with some acknowledgment of the question followed by the answer itself.
Preliminary experiments show that this consistent format generally improves model performance; we hypothesize that it assists the model in forming a chain of thought, similar to the ``let's think step-by-step'' prompt~\citep{kojima2022large,weichain}.
% Figure~\ref{fig:training_examples} presents 3 such examples.

We also include 13 training prompts with some degree of toxicity or malevolence.
We carefully write responses that partially or fully reject the command, and explain why the assistant will not comply.
There are also 30 prompts with similar issues in the test set, which we analyze in Section~\ref{sec:analysis}.

In addition to our manually authored examples, we sample 50 training examples from Super-Natural Instructions \citep{supernaturalinstructions}.
Specifically, we select 50 natural language generation tasks such as summarization, paraphrasing, and style transfer, and pick a single random example from each one.
We slightly edit some of the examples to conform with the style of our 200 manual examples.
While the distribution of potential user prompts is arguably different from the distribution of tasks in Super-Natural Instructions, our intuition is that this small sample adds diversity to the overall mix of training examples, and can potentially increase model robustness.

Manually creating diverse prompts and authoring rich responses in a uniform style is laborious.
While some recent works avoid manual labor via distillation and other automatic means~\citep{honovich2022unnatural, wang2022selfinstruct, alpaca, vicuna2023, sun2023principledriven}, optimizing for quantity over quality, this work explores the effects of investing in diversity and quality instead.